% Part: incompleteness
% Chapter: arithmetization-syntax
% Section: coding-formulas

\documentclass[../../../include/open-logic-section]{subfiles}

\begin{document}

\olfileid[pt]{inc}{art}{frm}
\olsection{Codificação de \printtoken{P}{formula}}

Depois de definir recursivamente de modo primitivo a relação
$\fn{Term}(x)$, podemos usá-la para definir da mesma maneira a relação
correspondente para !!{formula}s, $\fn{Frm}(x)$.

\begin{prop}
A relação $\fn{Atom}(x)$, que vale sse $x$ é o número de G\"odel de
!!a{formula} atômica, é recursiva primitiva.
\end{prop}

\begin{proof}
O número $x$ é o número de G\"odel de !!a{formula} atômica sse uma das
condições a seguir vale:
\begin{enumerate}
\item Existem $n$, $j < x$ e $z < x$ tais que, para cada $i < n$,
  $\fn{Term}((z)_i)$ e $x = $
\[
\Gn{\Obj P^n_j(} \concat \fn{flatten}(z) \concat \Gn{)}.
\]
\item Existem $z_1, z_2 < x$ tais que $\fn{Term}(z_1)$,
  $\fn{Term}(z_2)$ e $x = $
\[
\Gn{{\eq}(} \concat z_1 \concat \Gn{,} \concat z_2 \concat{\Gn{)}}.
\]
  \tagitem{prvFalse}{$x = \Gn{\lfalse}$.}{}
  \tagitem{prvTrue}{$x = \Gn{\ltrue}$.}{}
\end{enumerate}
\end{proof}

\begin{prop}
\ollabel{prop:frm-primrec}
A relação $\fn{Frm}(x)$, que vale sse $x$ é o número de G\"odel de
!!a{formula}, é recursiva primitiva.
\end{prop}

\begin{proof}
Uma sequência de símbolos~$s$ é !!a{formula} sse existe uma sequência
de formação de !!{formula}~$s_0$, \dots, $s_{k-1} = s$ que registra
como~$s$ foi formada a partir de !!{formula}s atômicas segundo as
regras de formação. O código de cada $s_i$ (e, na verdade, o código da
sequência $\tuple{s_0, \dots, s_{k-1}}$) é menor que o código~$x$
de~$s$.
\end{proof}

\begin{prob}
Apresente uma prova detalhada de
\olref[inc][art][frm]{prop:frm-primrec} seguindo as linhas da primeira
prova de \olref[inc][art][trm]{prop:term-primrec}.
\end{prob}

\begin{prop}
\ollabel{prop:freeocc-primrec}
A relação $\fn{FreeOcc}(x, z, i)$, que vale sse o $i$-ésimo símbolo da
fórmula com número de G\"odel~$x$ é uma ocorrência livre da variável
com número de G\"odel~$z$, é recursiva primitiva.
\end{prop}

\begin{proof}
Exercício.
\end{proof}

\begin{prob}
Prove \olref[inc][art][frm]{prop:freeocc-primrec}. Você pode usar o
fato de que toda subsequência de !!a{formula} que seja !!a{formula} é
uma sub-!!{formula} dela.
\end{prob}

\begin{prop}
A propriedade $\fn{Sent}(x)$, que vale sse $x$~é o número de G\"odel
de !!a{sentence}, é recursiva primitiva.
\end{prop}

\begin{proof}
Uma !!{sentence} é !!a{formula} sem ocorrências livres de
!!{variable}s. Portanto, $\fn{Sent}(x)$ vale sse
\begin{multline*}
\bforall{i<\len{x}}{\bforall{z<x}{}}\\
(\bexists{j<z}{z=\Gn{\Obj v_j}} \lif \lnot\fn{FreeOcc}(x,z,i)).
\end{multline*}
\end{proof}


\end{document}
